\chapter{Operators and their algebra}
\setwlchapter{algebra}
\section{What a declaration fixes}
The algebra of an operator is specified by its head, not its printed
letter. For fermions a mode is the tuple of indices \emph{after}
the type \fn{CR} or \fn{AN}. Changing any member of this tuple gives
an orthogonal mode. For spinful electrons the last index denotes spin;
an orbital, site, band or channel index belongs before it. Thus
$c^\dagger_{2a\uparrow}c_{2b\uparrow}$ transfers charge between
orbitals at fixed site and spin. SNEG keeps the same mode ordering
under conjugation and reverses the order of factors in a product.
\casestudy{spin-three-halves}{Changing the local spin multiplicity}

\section{Products, scalars and Hermitian conjugation}
With $\{c_i,c_j^\dagger\}=\delta_{ij}$, transposing neighbouring
factors generates a lower-degree term as well as a sign. Ordinary
Wolfram Language multiplication sorts its arguments and cannot encode
this operation. Use \fn{nc[a,b,c]} and \fn{pow[a,n]} for operator
products and their powers; use ordinary multiplication for commuting
coefficients. Declare real constants with \fn{snegrealconstants} and
complex amplitudes with \fn{snegcomplexconstants}. If a hopping
amplitude can carry flux, its conjugate is a different coefficient.
Real functions and real-valued indices are separate declarations:
$\epsilon(k)$ may be real even if its functional form is unknown.
\begin{wl}
snegfermionoperators[c];
snegcomplexconstants[v];
a = v nc[c[CR, 1, UP], c[AN, 2, UP]];
h = a + conj[a];
Expand[{h - conj[h], komutator[h, number[c[1]] + number[c[2]]]}]
\end{wl}
The two zeros answer different questions: the reverse hop is present
with the conjugate amplitude, and every term removes as many particles
as it creates. Hermitian pairing would pass the first test but generally
fail the second.

\section{Reference states}
In \fn{EMPTY} order an annihilator kills the vacuum. In \fn{SEA} order
the filled modes have negative energy or momentum index relative to
the chosen Fermi point. A creator below the surface therefore kills
the sea. At the Fermi point one must specify the state or resolve the
degeneracy; changing the order alone does not determine its physical
occupation. \fn{NONE} preserves products within a family and is useful
when explicitly comparing before and after a contraction. Changing
the ordering cannot undo simplifications of an already evaluated product.
\casestudy{sea}{A particle and a hole above the reference sea}

\section{Bosons, Majoranas and spin operators}
For a boson $[b,b^\dagger]=1$ replaces the fermion anticommutator.
A self-adjoint SNEG Majorana obeys $\{\gamma_i,\gamma_j\}=\delta_{ij}$:
its square is $1/2$, rather than one. A spin of size $s$, declared with
\fn{snegspinoperators}, acts in dimension $2s+1$ and commutes with
operators at other sites. Local products can be reduced using
$S^+S^-=s(s+1)-(S^z)^2+S^z$; a spin raising operator is not a
canonical boson. A Grassmann coefficient is yet another object:
it anticommutes with odd fermionic quantities but supplies no
creation--annihilation contraction.
\casestudy{pair-algebra}{Does a local fermion pair behave as a boson?}
\section{Transitions between many-body states}
Hubbard operators obey $X^{ab}X^{cd}=\delta_{bc}X^{ad}$ on the
same site. On different sites the parity of the transition matters:
moving an electron from one site to another still carries a Fermi
sign even if individual orbital operators have been eliminated.
\casestudy{graded}{Grading local transitions}
\begin{exercise}[Majorana normalization]\label{ex:majorana}
Write $f=(\gamma_1+i\gamma_2)/\sqrt2$. Determine $f^2$,
$\{f,f^\dagger\}$ and $f^\dagger f$ using the SNEG normalization.
\end{exercise}
